\chapter{Introduction}
\label{ch:introduction}
% ============================================================

\todo{Introduce the topic. Motivate the study. State RQ1 and RQ2 explicitly.
Give an overview of the thesis structure (one paragraph).}

\section{Motivation}
\paragraph{
Modern Large Language Models (LLMs) like ChatGPT, Claude and Gemini are able to produce high quality written text. Supplied with the right information, in theory these LLMs should be able to produce journalistic articles. As journalists are also trained and experienced in writing informative articles based on core information about the subject. As the speed of journalism has increased through the need of competition because of the connected world through ethernet, lowering the time spend on an article creates buesiness value in itself. The speed in which generative artificial intelegince (AI) can produce text is unmatched compared to humans. But can a general purpose AI generate the quality of text, for which customers are willing to buy and read an article? Are journalists able to implement the AI in their workflow to decrease their time spend on an article?}
\paragraph{
In this study I try to mimic the process behind journalisitic work for an application powered by AI. This tool is supplied to an cooperating editorial team. This team includes journalists with varying levels of experience and from different editorial departments to be able to obtain a representative perspective on journalistic work practices. Over a specified perios these journalists will include the provided tool in their day to day work and protocoll their work with it. Throught that I want to analyze the advantages and disadvantages working with generative AI in a real world scenario.}

\section{Research Questions}

The thesis investigates two core research questions:

\begin{description}
  \item[RQ1] To what extent does the use of generative AI tools affect the
    working efficiency and article quality of journalists?
  \item[RQ2] What measurable changes in writing style can be observed when
    comparing AI-assisted text production to unassisted writing?
\end{description}

\section{Thesis Structure}
Was reingehört:
Ein einleitender Satz, der den Gesamtaufbau ankündigt, gefolgt von einem Absatz (oder kurzen Sätzen) pro Kapitel – jeweils mit dem Kern-Inhalt, nicht nur dem Titel wiederholen.
Beispiel für deine Arbeit:

The remainder of this thesis is structured as follows. Chapter 2 provides the necessary background on generative AI, large language models, and stylometric analysis methods. Chapter 3 reviews related work on AI in journalism, human–AI collaboration, and computational writing-style analysis, and identifies the research gap this thesis addresses. Chapter 4 describes the study design, the developed web application, participant recruitment, data collection procedures, and the analysis methods applied. Chapter 5 presents the results, covering efficiency metrics, writing-style features, and qualitative interview findings. Chapter 6 discusses and interprets these findings, addresses limitations and threats to validity. Chapter 7 concludes the thesis by summarizing contributions, answering the research questions, and outlining directions for future work.